Les contraintes technologiques du projet sont les suivantes:
\begin{enumerate}
\item La première vient du fait que notre micro-controlleur
doit agir comme un host et une device usb en même temps. L'origine de la contrainte est
que le capteur optique utilisé est déjà interfacé en usb. Donc si nous voulons rajouter/
personnaliser des fonctionnalitées, il faut que le microcontrolleur lise les données du
capteur et les modifie avec les siennes avant de les transmettre à l'ordinateur toujours
sous forme de souris HID. Cela veut dire que l'on doit se restraindre à une device qui
possède 2 interfaces usb.

\item La deuxième provient du fait que pour scroll avec nos boutons il faut utiliser des io du
microcontrolleur qui lisent du capacitive touch et qui ne sont pas désactivées lorsque
ses deux interfaces usb sont utilisées.

\item Troisièmement il faut encore 2 io analogues pour le joystick qui ne soit pas bloqués
par les deux autres contraintes précédentes.
\end{enumerate}

Les contraintes de temps et de budget sont les suivantes:
\begin{enumerate}
\item Une première version de démonstration doit être disponible pour la fin de la saison
d'été du CCI.
\item Un budget en dessous de 20$ est nécessaire afin de garder le design accessible.
Par contre, la cible à viser est 10$.
\item La date butoir pour le design complet et finit de la souris est le 25 décembre 2024.
\begin{enumerate}
